\ExamPart{Câu tự luận}{3}{ }

\NQuestion Một tổ chuyên môn có $4$ thầy giáo, $5$ cô giáo và cần chọn một nhóm gồm đúng $1$ thầy, $1$ cô và $1$ thư ký trong số các người còn lại.
\begin{SubQuestions}
  \item Có bao nhiêu cách chọn nhóm nếu thư ký có thể là bất kì người nào chưa được chọn ở hai vị trí trước?
  \item Nếu thư ký bắt buộc là cô giáo, kết quả thay đổi như thế nào?
\end{SubQuestions}
\AnswerLine{$160$; khi thư ký là cô giáo thì có $80$ cách}
\SolutionBlock{\begin{SubQuestions}
\item Chọn $1$ thầy: $4$ cách. Chọn $1$ cô: $5$ cách. Khi đó còn lại $4+5-1=8$ người chưa được chọn để làm thư ký. Theo quy tắc nhân, số cách là
\[
4\cdot5\cdot8=160.
\]
\item Nếu thư ký bắt buộc là cô giáo thì sau khi chọn $1$ thầy và $1$ cô, còn lại $4$ cô giáo có thể làm thư ký. Số cách là
\[
4\cdot5\cdot4=80.
\]
\end{SubQuestions}}


\NQuestion Một khu đất hình chữ nhật có chiều dài hơn chiều rộng $4$ m. Sau khi người ta mở một lối đi rộng đều $1$ m ở phía trong quanh bốn cạnh thì phần đất còn lại ở giữa có diện tích $96$ m$^2$.
\begin{SubQuestions}
  \item Tính kích thước ban đầu của khu đất.
  \item Tính diện tích lối đi.
\end{SubQuestions}
\AnswerLine{$10\,\text{m}$ và $14\,\text{m}$; diện tích lối đi $44\,\text{m}^2$}
\SolutionBlock{Gọi chiều rộng ban đầu là $x$ m, chiều dài là $x+4$ m. Sau khi mở lối đi rộng $1$ m quanh bốn cạnh, phần đất còn lại có kích thước
\[
x-2 \quad \text{và} \quad x+2.
\]
Theo đề bài:
\[
(x-2)(x+2)=96 \Leftrightarrow x^2-4=96 \Leftrightarrow x^2=100.
\]
Vì $x>0$ nên $x=10$. Vậy khu đất ban đầu có kích thước $10$ m và $14$ m.
Diện tích ban đầu là $10\cdot14=140$ m$^2$, diện tích phần còn lại là $96$ m$^2$, nên diện tích lối đi là $44$ m$^2$.}

\NQuestion Trong mặt phẳng tọa độ $Oxy$, cho ba điểm $A(1;1)$, $B(5;3)$, $C(3;-1)$.
\begin{SubQuestions}
  \item Tính tọa độ các vectơ $\overrightarrow{AB}$ và $\overrightarrow{AC}$.
  \item Viết phương trình đường thẳng $BC$.
  \item Tính khoảng cách từ điểm $A$ đến đường thẳng $BC$.
\end{SubQuestions}
\AnswerLine{$\overrightarrow{AB}=(4;2),\ \overrightarrow{AC}=(2;-2),\ BC:2x-y-7=0,\ d(A,BC)=\dfrac{6}{\sqrt5}=\dfrac{6\sqrt5}{5}$}
\SolutionBlock{\begin{SubQuestions}
\item Ta có
\[
\overrightarrow{AB}=(5-1;3-1)=(4;2),\qquad \overrightarrow{AC}=(3-1;-1-1)=(2;-2).
\]
\item Đường thẳng $BC$ qua $B(5;3)$ và $C(3;-1)$ có hệ số góc
\[
m=\frac{-1-3}{3-5}=2.
\]
Do đó phương trình là
\[
y-3=2(x-5) \Leftrightarrow 2x-y-7=0.
\]
\item Khoảng cách từ $A(1;1)$ đến $BC$ là
\[
d(A,BC)=\frac{|2\cdot1-1-7|}{\sqrt{2^2+(-1)^2}}=\frac{6}{\sqrt5}=\frac{6\sqrt5}{5}.
\]
\end{SubQuestions}}


\NQuestion (Câu khó) Trong mặt phẳng tọa độ $Oxy$, cho hai điểm $A(-1;2)$, $B(3;4)$ và điểm $M$ thuộc trục tung $Oy$.
\begin{SubQuestions}
  \item Đặt $M(0;t)$, hãy biểu diễn $MA^2+MB^2$ theo $t$.
  \item Tìm tọa độ điểm $M$ để tổng $MA^2+MB^2$ nhỏ nhất.
  \item Tính giá trị nhỏ nhất đó.
\end{SubQuestions}
\AnswerLine{$M(0;3),\ \min(MA^2+MB^2)=12$}
\SolutionBlock{\begin{SubQuestions}
\item Với $M(0;t)$, ta có
\[
MA^2=(0+1)^2+(t-2)^2=1+(t-2)^2,
\]
\[
MB^2=(0-3)^2+(t-4)^2=9+(t-4)^2.
\]
Do đó
\[
MA^2+MB^2=10+(t-2)^2+(t-4)^2.
\]
Khai triển:
\[
MA^2+MB^2=2t^2-12t+30=2(t-3)^2+12.
\]
\item Biểu thức đạt giá trị nhỏ nhất khi $t=3$. Vậy $M(0;3)$.
\item Giá trị nhỏ nhất là $12$.
\end{SubQuestions}}
